\documentclass[../numerical_methods.tex]{subfiles}
\begin{document}
\section{Aula 04 - 12 de Março, 2026}
\subsection{Motivações}
\begin{itemize}
	\item Introdução ao Numpy.
\end{itemize}
\subsection{Numpy}
O \texttt{numpy} é um pacote no Python para a computação orientada às matrizes, conhecida como \textbf{ processamento matricial}; nele, as classes e métodos são majoritariamente escritos em C, garantindo eficiência computacional, e possui muitos recursos para computação científica, especialmente para resolver problemas de álgebra linear.
Ele é um pacote muito utilizado, ao ponto de vários outros, como \texttt{pandas}, \texttt{matplotlib} (que veremos depois) e \texttt{sklearn}, dependerem dele.

Para começar a usar o numpy, não basta apenas ter o Python, ele é um pacote, e precisa ser importado após baixado; supondo que tenha instalado, pode importá-lo como ``\texttt{import numpy as np}'', para que toda chamada do pacote seja com np na frente, ou ``\texttt{import numpy}'' seguido de ``from numpy import *''.
A partir disso, a estrutura de dados básica do numpy são os chamados \textit{arrays}, que funcionam similarmente às listas em Python, mas com alguns detalhes:
\begin{itemize}
	\item Todo elemento em um array precisa ser do mesmo tipo, normalmente um float ou um int;
	\item Cada dimensão de um array é chamado de eixo/\textit{axis};
	\item Os eixos são enumerados a partir de 0; e
	\item Os elementos são acessados de forma parecida com as listas, utilizando colchetes [].
\end{itemize}
Os arrays são a estrutura que viabilizam a realização eficiente de operações numéricas envolvendo quantidades enormes de dados, e são uma forma mais eficiente das listas que conhecemos, tanto é que uma das formas de construí-los é convertendo de listas para arrays; outras, seriam os próprios métodos do numpy ou a conversão de \texttt{DataFrames}:
\begin{listing}
	\caption{Criando arrays a partir de listas.}
	\begin{minted}[bgcolor = DarkBackground, linenos=true]{Python}
 import numpy as np 

 # Criando um array unidimensional a partir de uma lista 

 lista = [1,3,5,7,9,10]
 array = np.array(lista)
 print(array)
 print(lista)
 # Note que, na impressão de um array, as vírgulas não aparecem 
 # Mas, nas listas, elas aparecem.
 
 # Criando um array bidimensional

 matriz = [[1,3,5,7,9,11], [2,4,6,8,10,12], [0,1,2,3,4,5]] 
 array_bidim = np.array(matriz)
 print(array_bidim)

 \end{minted}

	\caption{Criando arrays pelo numpy.}
	\begin{minted}[bgcolor = DarkBackground, linenos=true]{Python}
  # Criar uma array a partir dos métodos numpy é parecido com criar 
  # listas de template. Algumas formas: 

  # Arrays unidimensionais 

  array_zeros = np.zeros((8)) # array com oito elementos iguais a 0
  array_uns = np.ones((10)) # array com dez elementos iguais a 1 
  array_cont = np.arrange(10) # array com elementos de 0 a 9
  array_part = np.linspace(1,2,5) # array com 5 elementos igualmente 
  # espaçados no intervalo de 1 até 2.

  # Arrays bidimensonais 

  cinco_tres = np.zeros((5,3)) # array 5x3 de zeros
  cinco_dez = np.ones((5, 10)) # array com 5x10 de uns
  identidade = np.identity(3) # matriz identidade 3x3
 \end{minted}
\end{listing}

Os objetos correspondentes às arrays do numpy, entendidas pelo computador como objetos do tipo \texttt{ndarrays}, possuem diversos atributos que podem ser acessados:
\begin{itemize}
	\item \textbf{ndarray.ndim:} número de eixos/dimensões da array;
	\item \textbf{ndarray.shape:} tupla de inteiros indicando o tamanho da array em cada dimensão;
	\item \textbf{ndarray.size:} número total de elementos da array;
	\item \textbf{ndarray.dtype:} tipo dos elementos da array;
	\item \textbf{ndarray.itemsize:} tamanho, em bytes, de cada elemento da array;
	\item \textbf{ndarray.data:} buffer de memória contendo os elementos do array.
\end{itemize}
Vale observar que arrays unidimensionais podem ter uma das dimensões livres: um array com \textit{shape} (5,) possui apenas 1 eixo, enquanto que um com (5,1) possui 2.

Um conceito importante de ser entendido ao manipular arrays é o de \textit{view}, que é criada em algumas ocasiões, como ao realizar um \textit{slicing}, que será visto adiante, e faz referência a uma parte de uma array; é importante manter em mente que alterar os elementos de um view irá alterar os elementos da array original, então caso seja necessário preservá-la, deve-se fazer uma cópia usando o método \texttt{copy}. 

Os índices de uma array variam de 0 até \(k_{i}-1\), sendo \(k_{i}\) o número de elementos na i-ésima dimensão, e a iteração é feita linha pro linha; assim, se o desejo for acessar os elementos, deve-se utilizar um laço duplo, ou a versão \textit{flat} da array:
\begin{minted}[bgcolor = DarkBackground, linenos=true]{Python}
import numpy as np

matriz = [[1,3,5,7,9,11], [2,4,6,8,10,12], [0,1,2,3,4,5]] 
array_bidim = np.array(matriz)

# percorrendo as linhas 
for indice, linha in enumerate(array_bidim):
    print("linha", indice, " = ", linha)

# percorrendo elementos 
for indice, linha in enumerate(array_bidim):
    for coluna, elemento in enumerate(linha):
        print("Elemento na posição ", i, ' ', j, ' = ', elemento)

# o atributo 'flat' coloca todos os elementos em uma lista 
for elemento in array_bidim.flat:
    print(elemento)

# pode-se utilizar os índices diretamente 
for linha in range(array_bidim.shape[0]): # numero linhas
    for coluna in range(array_bidim[1]): # numero colunas
        print('Elemento', i, ' ', j, ' = ', array_bidim[i, j])
\end{minted}
O problema de usar os laços é que eles são pouco eficazes, computacionalmente, e existe uma alternativa conhecida como \textit{slicing} (podendo ser usado também para listas), basicamente uma visão do array original (gera uma view), sem que o dado seja copiado. Ao utilizar o glicina, tenha em mente que índices omitidos correspondem a enxergar toda a dimensão omitida.

\begin{minted}[bgcolor = DarkBackground, linenos=true]{Python}
matriz = [[1,3,5,7,9,11], [2,4,6,8,10,12], [0,1,2,3,4,5]] 
array_bidim = np.array(matriz)

# slicing de colunas 
## array[indice_linha, :] recupera a linha inteira do indice 
## indicado. 
print(array_bidim[1,:]) # recupera toda a linha 1.

# slicing de linhas 
## array[:, indice_coluna] recupera a coluna inteira do indice 
## indicado 
print(array_bidim[:, 2]) # recupera toda a coluna 2. 

# slicing em blocos 
## podemos estabelecer a recuperação em linhas e em colunas usando 
## blocos. 
print('Linhas 0 e 2: ', array_bidim[[0, 2]])
print('Colunas 0, 2 e 5: ', array_bidim[:[0,2,5]])
print('Elementos nas posições (0,0), (1,2) e (2,5): ',
array_bidim[[0,1,2], [0,2,5]])

''' 
quando duas listas, uma com os indices de linha, como [0,1,2], e 
outra com indices de colunas, como [0,2,5], são passadas para o 
slicing, os elementos são recuperados em pares! 
Assim, as duas listas devem ter o mesmo número de elementos/indices.
'''

# linha ou coluna fixada 
## neste caso, o slicing gera um array unidimensional 
print('[:2, 1]: ', array_bidim[:2, 1].shape)
print('\n[:, -1]: ', array_bidim[:, -1].shape)
\end{minted}

Também dá para acessar os elementos usando alguma forma de condicional, um valor booleano permitindo fazer as atribuições de maneira mais eficiente e elegante, conhecida como \textbf{máscara booleana}:
\begin{minted}[bgcolor = DarkBackground, linenos=true]{Python}
matriz = [[1,3,5,7,9,11], [2,4,6,8,10,12], [0,1,2,3,4,5]] 
array_bidim = np.array(matriz)

# Mostra os valores-verdade de cada entrada da matriz
mascara = (array_bidim > 7)
print(mascara)

# Mostra apenas os elementos que satisfazem a condição
print(array_bidim[mascara])

# Modifica apenas os elementos que satisfazem a condição
array_bidim[mascara] = 0 
print(array_bidim)
\end{minted}

\subsection{Manipulando Arrays}
O numpy traz em si alguns métodos para modificar a estrutura das arrays. Dentre eles, temos
\begin{itemize}
	\item \textbf{\underline{reshape}:} permite reformatar (reshape) o array, modificando o número de linhas e colunas delas; porém, o novo formato (a nova shape) deve ter o mesmo número de elementos do array original: uma array com 9 elementos em linhas (1x9) pode virar uma array 3x3 (9 elementos);
	\item \textbf{\underline{ravel}:} concatena as linhas da matriz em uma array unidimensional e gera uma view, tal que, se um elemento for modificado, a array original também será modificada;
	\item \textbf{\underline{flatten}:} concatena as linhas da matriz em uma array unidimensional, mas faz uma cópia dos elementos;
	\item \textbf{\underline{.T}:} transpõe a array (e matrizes!);
\end{itemize}

\begin{minted}[bgcolor = DarkBackground, linenos=true]{Python}
array_unidim = np.arange(9)

# reformata de um array aleatório com números entre 0 e 9 
# para uma matriz 3x3
reformatar = array_unidim.reshape((3,3))

# concatena as linhas da array e cria uma cópia 
concatena_flatten = array_unidim.flatten()
print('array concatenada com flatten: ' concatena_flatten)

# podemos alterar a nova sem alterar a original
concatena_flatten[2]=-1
print('array concatenada e modificada: ' concatena_flatten)
print(array_unidim)

# mas não é o caso com o ravel
concatena_ravel = array_unidim.ravel()
print('array concatenada com ravel: ', concatena_ravel)

concatena_ravel[2]=-1 
print('array concatenada com ravel e modificada: ', concatena_ravel)
print('array original: ', array_unidim)
\end{minted}

\subsection{I/O: Input e Output com Numpy}
É possível utilizar métodos para ler e escrever arrays em arquivos de textos novo. Sobre o carregamento de arquivos, os valores no arquivo devem estar separados por espaços, e cada linha dele vira uma linha da array resultante, sendo que vários parâmetros podem ser especificados com o método \texttt{loadtxt}; para os métodos de escrita de arquivos novos, também existem parâmetros que podem ser especificados. Leiam as documentações!
\begin{minted}[bgcolor = DarkBackground, linenos=true]{Python}
# Gerando 18 números aleatórios entre 0 e 10 e reformata em 
# uma matriz 3x6; depois, salva ela em um arquivo chamado 
# exemplo_escrita.txt

import numpy as np

array = np.random.randint(0,10,18).reshape(3,6)
np.savetext('exemplo_escrita.txt', x, fmt = '%d')

# Agora, carregamos o arquivo salvo, forçando que os elementos 
# da matriz sejam inteiros. 
array_arquivo = np.loadtxt('exemplo_escrita.txt', dtype=int)
print(array_arquivo)

# veja o que acontece quando `dtype` não é especificado! 
array_teste = np.loadtxt('exemplo_escrita.txt')
print(array_teste)
# Elementos da matriz são do tipo "float".
\end{minted}
\end{document}
